\section{Introduction}
A whole-line operator can have only three spectral values and still acquire
rich spectral structure when restricted to a long interval. The sine
Loewner operator provides an explicit example. Its whole-line Fourier
representation exchanges two adjacent frequency windows, with spectrum
contained in $\{-1,0,1\}$. Spatial truncation produces a logarithmic
transition near the two outer spectral bands and isolated states in the
central gaps. The outer transition admits a uniform trace law; the gap
states follow a one-period boundary flow and interact across the interval
at an explicitly determined algebraic rate.

The finite matrices are
\begin{equation}\label{eq:2.1}
 A_N(\varepsilon)_{mn}=
 \begin{cases}
 \dfrac{\sin(2\pi m\varepsilon)-\sin(2\pi n\varepsilon)}{\pi(m-n)},&m\ne n,\\[1ex]
 2\varepsilon\cos(2\pi m\varepsilon),&m=n,
 \end{cases}\qquad {-N}\le m,n\le N,
\end{equation}
where $N\ge1$, $d=2N+1$ and $0<\varepsilon<1/2$. Their critical
integral operator on $L^2[-1,1]$ has kernel
\begin{equation}\label{eq:2.3}
 K_\tau(x,y)=\frac{\sin(2\pi\tau x)-\sin(2\pi\tau y)}{\pi(x-y)},
 \qquad K_\tau(x,x)=2\tau\cos(2\pi\tau x).
\end{equation}
Both are self-adjoint contractions and commute with reflection.
The kernel is the commutator of multiplication by $\sin(2\pi\tau x)$
and the finite Hilbert transform with kernel $1/[\pi(x-y)]$.
The critical scale is $N\varepsilon$; an exact covariance comparison
uses the slightly different parameter $d\varepsilon/2$.

\subsection*{Three main conclusions}
\emph{A uniform trace law and a finite-matrix bridge.}
For $r_t=\sqrt{t^2+(1-t)^2}$, set
\begin{equation}\label{T10:eq:kappa}
 \varkappa(f)=\frac1{2\pi^2}\int_0^1
 \frac{f(r_t)+f(-r_t)-f(1)-f(-1)}{t(1-t)}\,dt.
\end{equation}
The integral converges for H\"older $f$: the endpoint numerator is
$O([t(1-t)]^\alpha)$ for its H\"older exponent $\alpha>0$.
\begin{theorem}[Uniform trace law]\label{T10:conj:law}
Every real $f\in C^2([-1,1])$ with $f(0)=0$ satisfies, for $\tau\ge2$,
\begin{equation}\label{T10:eq:law}
 \tr f(\mathcal K_\tau)
 =2\tau[f(1)+f(-1)]+\varkappa(f)\log\tau+\mathsf R_\tau(f),
\end{equation}
where, with the absolute constant $C_{\mathrm{tr}}=143$,
\begin{equation}\label{rev5:eq:uniform-remainder}
 |\mathsf R_\tau(f)|\le C_{\mathrm{tr}}
                    (\|f'\|_\infty+\|f''\|_\infty).
\end{equation}
The bound is uniform in $f$ and $\tau$.
\end{theorem}
The two bulk masses lie at $\pm1$. The logarithmic correction is
supported on $g\le|\lambda|\le1$, where $g=1/\sqrt2$, with the
endpoint subtraction in \eqref{T10:eq:kappa} essential. The theorem
extends to globally Lipschitz tests that are $C^{1,1}$ only on these
outer bands. It gives all Schatten powers $|x|^p$, $p\ge1$, closed
determinant coefficients and signed transition counts. In particular,
for $w\ge0$ the coefficient in $\log\det(I+w\mathcal K_\tau^2)$ is
\[
 -\frac4{\pi^2}\arcsin^2\!\left(\frac12\sqrt{\frac{2w}{1+w}}\right),
\]
which equals $-1/9$ at $w=1$.

Theorem~\ref{v11:thm:direct-transfer} compares $A_N(\varepsilon)$
with $\mathcal K_{d\varepsilon/2}$ independently of dimension.
After the exact linear correction, its $C^2$ trace error is at most
$L_\varepsilon\|f''\|_\infty$, with
$L_\varepsilon=2\varepsilon\csc(2\pi\varepsilon)-1/\pi
=O(\varepsilon^2)$. The Lipschitz error is $O(\varepsilon)$.
Thus the trace laws also apply at fixed sampling scale, uniformly
when $\varepsilon$ stays below $1/2-\delta$, $\delta>0$.

\emph{Physical gap states and one-period flow.}
Dilation identifies $\mathcal K_\tau$ with the compression to
$(-\tau,\tau)$ of the fixed whole-line exchange $\mathcal G$ with
kernel $[\sin(2\pi X)-\sin(2\pi Y)]/[\pi(X-Y)]$.
Let $B_b=P_{(b,\infty)}\mathcal G P_{(b,\infty)}$.
For each fixed $\alpha\in(-g,0)\cup(0,g)$ there is exactly one
$b$ modulo one for which $\alpha$ is an eigenvalue of $B_b$, and
that eigenvalue is simple. Long-interval clusters converge to the
two boundary spectra with exact multiplicity
(Theorems~\ref{v11:thm:gap-limits} and
\ref{boundary:thm:one-cell-flow}). This information is finer than the
bulk and logarithmic trace terms.

\emph{The full algebraic interaction and a certified nonzero example.}
For an isolated level $\alpha$ shared by the two boundaries, let
$\eta_L,\eta_R$ be the normalized physical tail amplitudes of
Section~\ref{unified:sec:interaction}. For an interval of length $L$
with fixed endpoint phases, the two nearby eigenvalues are
\[
 \alpha+L^{-\gamma(\alpha)}\lambda_j(M_\alpha)
       +o(L^{-\gamma(\alpha)}),\qquad j=1,2,
\]
where
\[
 \gamma(\alpha)=\frac2\pi\arctan\sqrt{1-2\alpha^2},\qquad
 M_\alpha=-\frac{2\pi(1-\alpha^2)\sqrt{1-2\alpha^2}}{\alpha}
 \begin{pmatrix}0&\overline{\eta_L}\eta_R\\
                 \overline{\eta_R}\eta_L&0\end{pmatrix}.
\]
This matrix includes the complementary-resolvent correction, which
has the same order as the direct overlap
(Theorem~\ref{v13:thm:leading-interaction}). Centered intervals have
opposite even and odd coefficients; the odd state lies farther from
zero when the amplitude is nonzero. Exact interval certificates
isolate a phase-zero mode in $(0.405,0.428)$, prove its uniqueness in
$[0.3565,0.4765]$, and establish $|\eta|>1/100$ for every normalized
physical eigenfunction there. They therefore give an unconditional
instance of the nonzero splitting law. General nonvanishing is not
assumed or proved.

\subsection*{Other consequences and scope}
At small continuum scale, the four leading signed/parity branches begin
at orders $\tau,\tau^3,\tau^5,\tau^7$. For a simple isolated parity
eigenvalue in the finite critical limit, physical eigenvector
reconstructions converge weakly to zero while magnification reveals
nontrivial endpoint profiles. These are distinct from the long-window
boundary modes.
Physical boundary localization also gives continuous phase-dependent third terms.
For finite matrices these hold whenever $\varepsilon_N\to0$ and
$\tau_N=(2N+1)\varepsilon_N/2\to\infty$, without a further condition
on $(2N+1)\varepsilon_N^3$ or convergence of the fractional phases.
Reflection resolves information invisible in the leading signed trace:
for $w>0$,
\[
 \frac{\det(I+w\mathcal K_{\tau,+}^2)}
      {\det(I+w\mathcal K_{\tau,-}^2)}\longrightarrow1+\frac w2,
\]
where $+$ and $-$ denote spatial parity. For a fixed finite
commensurate frequency family, a bipartite bulk graph is sign-symmetric,
whereas additive frequency triples produce an odd logarithmic term.

The estimates have different scopes, summarized below. In particular,
relative gap indices do not determine unspecified absolute ranks, and
no uniform interaction law at zero or at a band edge is asserted.
\begin{center}
\begin{tabular}{@{}p{.28\textwidth}p{.65\textwidth}@{}}
\hline
Question & Parameter and regularity range\\\hline
Critical finite limit & $N\varepsilon_N\to\tau\in(0,\infty)$;
all fixed signed/parity eigenvalue branches.\\
Uniform trace remainder & $\tau\ge2$; global $C^2$, or global Lipschitz
and band-local $C^{1,1}$.\\
Finite covariance transfer & $0<\varepsilon<1/2$; constants independent
of $N$, with stated dependence on $\varepsilon$.\\
Phase-dependent constant & Global Lipschitz and band-local $C^2$;
continuous periodic phase profile.\\
Boundary interaction & Fixed isolated $\alpha\in(-g,0)\cup(0,g)$;
fixed endpoint phases and integer increases of interval length.\\
Certified nonvanishing & The isolated phase-zero mode enclosed in
$(0.405,0.428)$.\\\hline
\end{tabular}
\end{center}

\subsection*{Relation to earlier work}
Loewner matrices are classical \cite{Loewner1934}; their inertia is
studied in \cite{BhatiaFriedlandJain2016}.  De Andrade Silva recorded the arithmetic
divided-difference block \cite{SilvaPrime,SilvaFramework}; its finite
arithmetic dictionary and cutoff-derivative geometry occur in
\cite{GroskinDictionary,GroskinPrimeEvent}.  Compressed translations
appear in \cite[Section~5.1]{GuoPoleFree}, and the commutator framework
is classical \cite{Calderon1965}.
Positive time--band concentration and discrete prolate sequences are
classical \cite{Fuchs1964,LandauPollak1962,LandauWidom1980,Slepian1965,Slepian1978,SlepianPollak1961};
quantitative bounds include \cite{BonamiJamingKaroui2021,Karnik2021}.
No priority is claimed for these objects or for the variable $N\varepsilon$.
The same matrix also appears in Shi
\cite[eqs.~(3.4)--(3.7)]{Shi2026}: in that paper's notation,
$2L\,\mathbf G_{N,L}(x)=A_N(1-x/L)$ for $L/2<x<L$, including the
diagonal entries.  Shi studies a finite arithmetic model and its
rank-two displacement structure; the critical scaling and uniform
trace asymptotics considered here are different conclusions.

The two-term structure of Theorem~\ref{T10:conj:law} is that of the
Widom--Sobolev formula \cite{BollmannMueller2024,Sobolev2013,Widom1982}.
In the notation of \cite[eq.~(2.27)]{BollmannMueller2024},
$\varkappa(f)=2\,\mathfrak U(f;X_{1,0},X_{0,-1})$, where $X_{1,0}$ and
$X_{0,-1}$ are the two one-sided limits of the Bloch symbol of
Proposition~\ref{T4:prop:bloch}, the factor $2$ counts the two ends of
the truncation, and $\pm\sqrt{t^2+(1-t)^2}$ is the nonzero spectrum of
their linear interpolation (Section~\ref{v6:sec:density}).  Matrix-valued trace formulas, including the jump integral displayed in
\cite[eq.~(1.3)]{Shao2003}, provide direct precedents for the reduced
three-dimensional symbol. Block determinant asymptotics and
operator-valued Szeg\H{o}--Widom theory supply further settings
\cite{BasorEhrhardtVirtanen2024,BoettcherSilbermann1994};
Bloch-function trace asymptotics for a lattice-periodic operator appear in
\cite{PfirschSobolev2018}. The infinite physical fiber must be reduced
with control of the compression error before the finite-matrix reference
can be used. The positive logarithmic counting and
logistic rank laws are already in
Landau--Widom~\cite[pp.~469--470]{LandauWidom1980}.
What is proved here is the reduction of the exchange symbol to a fixed
three-dimensional jump, the resulting trace law with a remainder bounded
by the first two derivative norms uniformly in $f$ and $\tau$, its
band-local extension, the folded signed profile, and the quantified
dimension-independent transfer back to finite matrices.  The Hilbert--Carleman comparison has a direct antecedent in
Otte~\cite{Otte2024}, and the smooth compression estimate is a version of
Laptev--Safarov~\cite{LaptevSafarov1996}. We prove all operator comparisons
needed here. Recent work also treats moving thresholds and low-regularity traces for
positive localization operators \cite{Azimifard2026b,KulikovDamLarsen2026}.
Our signed transition estimate is uniform on compact subsets of the outer
bands; it does not identify individual eigenvalues with positive prolate
eigenvalues or assert their sharper counting errors for this operator.



One-period edge flow has precedents for periodic differential
operators \cite[Proposition~3.21]{Gontier2020}; the proof here uses
an analytic bandlimited covariance and a relative projection index.
The endpoint model is related to classical finite-Hilbert-transform
analysis \cite{KoppelmanPincus1959,BertolaEtAl2020,AlAifariKatsevich2014}.
Our frequency pencil retains the same-half Hilbert blocks together
with a phase-dependent exchange; the scalar reductions for adjacent
or overlapping intervals do not directly diagonalize that pencil.
The full complementary correction has the structure of a
Feshbach--Schur reduction \cite{DussonSigalStamm2021}, but its
coefficient is evaluated here for the physical boundary modes.
Block Toeplitz determinant constants
\cite{BasorEhrhardtVirtanen2024} and scalar
Toeplitz-plus-Hankel parity formulas \cite{BasorEhrhardt2002}
provide the reference results; we prove the localization and reflected
limits needed to transfer them to the present operator-valued symbol.



\subsection*{Organization}
Section~\ref{unified:sec:foundations} gives the exchange representation,
strict signed monotonicity and critical limit.
Section~\ref{unified:sec:trace} develops the logarithmic density,
coefficients and complete uniform-remainder proof.
Section~\ref{unified:sec:finite} gives finite-scale estimates and
endpoint geometry; Section~\ref{unified:sec:frequency} treats
commensurate frequencies.
The boundary analysis proceeds from phase constants
(Section~\ref{unified:sec:phase}) to half-line flow and the endpoint
equation (Section~\ref{unified:sec:gap}), the full interaction
(Section~\ref{unified:sec:interaction}) and its certificates
(Section~\ref{unified:sec:certificate}).
Section~\ref{unified:sec:parity} resolves parity traces and counts.
The endpoint spectral measure and differential obstruction in
Section~\ref{unified:sec:gap}, together with the final open questions,
delineate what these methods do and do not determine. All mathematical proofs are
included in the paper; the ancillary files contain reproducible
programs and exact rational witnesses.
